%! End Behavior and Asymptotes — Unit 2, Lesson 2.5 — Student Packet Cover Sheet
\documentclass[10pt]{article}
\usepackage{atda-article}
\usepackage{atda-boxes}
\usepackage{ltablex}
\keepXColumns

\begin{document}

% Full-bleed royal banner
\begin{tikzpicture}[remember picture, overlay]
  \fill[royal] (current page.north west)
    rectangle ([yshift=-0.9in]current page.north east);
\end{tikzpicture}
\vspace{-0.8in}

\begin{center}
  {\color{white}\textbf{\LARGE Algebra, Trigonometry, and Data Analysis}}\\[4pt]
  {\color{mistmid}\large\textbf{Unit 2: Functions, Transformations, and Graph Analysis}}\\[6pt]
  {\color{white}\textbf{\large Lesson 2.5 \quad End Behavior and Asymptotes}}
\end{center}

\vspace{0.22in}
\namedateperiod
\vspace{0.18in}

\begin{learningtargetbox}
\small
\begin{enumerate}[leftmargin=1.5em, itemsep=5pt]
  \item \ldots{} describe the \textbf{end behavior} of a graph --- what the \emph{outputs} do as the
        \emph{inputs} run off to the far left and the far right --- using one of three answers: they
        climb without bound, they fall without bound, or they level off toward a number.
  \item \ldots{} identify a \textbf{horizontal asymptote} from a graph and name it as an
        \emph{equation} of a line, not as a bare number.
  \item \ldots{} identify a \textbf{vertical asymptote} from a graph, connect it to the input that is
        missing from the domain, and explain why the graph can never cross it.
  \item \ldots{} explain why ``gets closer and closer'' is not ``arrives,'' and use that to say why a
        graph may have no zeros and no absolute maximum or minimum.
  \item \ldots{} interpret both kinds of asymptote in the situation the graph models --- the level a
        quantity settles toward, and the input a model is not allowed to take.
\end{enumerate}
\end{learningtargetbox}

\vspace{0.14in}

\begin{tocbox}
\small
\renewcommand{\arraystretch}{1.5}
\begin{tabularx}{\linewidth}{@{} c l X r @{}}
\rowcolor{mist}
\textbf{\#} & \textbf{Component} & \textbf{Description} & \textbf{Score} \\
\midrule
1 & Warm-Up        & Halving forever, and a number that is never negative & \blank{1.2cm} \\
2 & Guided Notes   & What a graph does at the far ends, and the two kinds of line it approaches but
                     never reaches                                 & \blank{1.2cm} \\
3 & Group Activity & A drug leaving the bloodstream, a typing speed with a ceiling, and an asymptote
                     that moved                                    & \blank{1.2cm} \\
4 & Exit Ticket    & One warming soda, and a classmate who put the graph on the line
                                                                   & \blank{1.2cm} \\
5 & Homework       & A refilling reservoir, a graph with a hole in its domain, and a look ahead
                                                                   & \blank{1.2cm} \\
\midrule
  & \multicolumn{2}{r}{\textbf{Total}} & \blank{1.2cm} \\
\end{tabularx}
\end{tocbox}

\vspace{0.14in}

\begin{remindbox}
\small
This lesson covers \texttt{AFDA.AF.2f} (describe the end behavior of a function given its graph) and
\texttt{AFDA.AF.2g} (identify horizontal and/or vertical asymptotes from a graph, if they exist).
Everything is still read \textbf{off a graph}. One sentence runs the whole packet: \textbf{getting
closer and closer is not the same as arriving.} Lesson 2.4 taught you to say which axis a number came
from; today you also have to say whether the graph \emph{reaches} a value or only \emph{approaches}
it --- and an asymptote is a \textbf{line}, so it gets an equation, never a bare number.
\end{remindbox}

\end{document}
